\section{Introduction}
\label{sec:intro}

Can large language models only express certain personas after being shown hundreds or thousands of examples of that persona in action? If so, short system prompts would access only a fraction of the model's full representational capacity, and the ``true'' persona space would be much larger than what few-shot prompting can reach. Conversely, if \emph{every} persona is accessible with minimal context, the practical persona space is bounded by what can be concisely described---vast, but shallow.

{\bf Why this matters.}
LLMs are increasingly deployed to simulate diverse perspectives: political viewpoints for deliberation platforms \citep{tencent2024personahub}, personality profiles for social science \citep{wang2025geometry}, and character roles for interactive fiction. The design of these systems depends critically on whether a short prompt suffices to activate any target persona or whether some personas require extensive contextual conditioning. If context-heavy personas exist, retrieval-augmented persona systems become necessary; if they do not, simple prompting strategies are sufficient.

{\bf The gap.}
Prior work on persona elicitation has tested only narrow ranges of in-context example counts. \citet{choi2024picle} demonstrated dramatic improvements from $K{=}0$ to $K{=}3$ in-context examples on Llama-2-7B, but never tested beyond $K{=}3$. Long-context in-context learning studies \citep{bertsch2024longcontext} scaled to 2000+ demonstrations but only for classification tasks, not persona elicitation. \citet{chen2025personavectors} showed that many-shot prompting progressively shifts persona vector projections, but tested only 0--20 examples. No study has systematically measured persona elicitation across a wide range of context lengths ($K{=}0$ to $K{=}200$+) to determine whether ``hard'' personas exist that require extensive conditioning.

{\bf Our approach.}
We conduct the first systematic study of persona elicitation accuracy as a function of in-context demonstration count $K$, spanning $K{=}0$ to $K{=}200$, across 25 diverse personas from the Anthropic Model-Written Evals benchmark. We test \gptmini with two complementary experiments: (1) a binary agreement task measuring action consistency across nine $K$ values, and (2) an open-ended generation task measuring persona distinctiveness via LLM-as-judge scoring.

{\bf Preview of results.}
We find no evidence for context-heavy personas. Every tested persona achieves $\geq$95\% accuracy at $K{=}0$, with a mean improvement of less than 0.4\% from $K{=}0$ to $K{=}200$. Generation distinctiveness shows a marginally significant but tiny improvement (+0.18 on a 10-point scale, $p{=}0.038$). The only personas exhibiting any difficulty at $K{=}0$ are AI-risk-related behaviors (\eg self-replication, sandbox escape), where safety training creates initial resistance that a handful of examples can overcome.

Our contributions are:
\begin{itemize}[leftmargin=*,itemsep=0pt,topsep=0pt]
    \item We conduct the first systematic study of persona elicitation accuracy across a wide range of in-context demonstration counts ($K{=}0$ to $K{=}200$), finding that all 25 tested personas saturate at or near $K{=}0$ on a frontier model.
    \item We propose the \emph{shallow persona hypothesis}: for sufficiently capable models, the accessible persona space is bounded by description length rather than context length, implying that all practically relevant personas can be specified with short prompts.
    \item We complement binary persona agreement with an open-ended generation distinctiveness evaluation, finding only marginal gains from additional context (+0.18/10 mean improvement).
    \item We identify safety-trained AI-risk personas as the only category showing any context sensitivity, attributing this to RLHF guardrail resistance rather than genuine persona depth.
\end{itemize}
